\chapter{Proiectare de detaliu și implementare}\label{ch:proiectare}
\pagestyle{fancy}

Acest capitol prezintă implementarea soluției fundamentate în capitolul~\ref{ch:analysis}. Sunt descrise organizarea modulară a codului, structurile de date utilizate în fiecare etapă, algoritmii implementați și principalele decizii determinate de particularitățile datelor și ale bibliotecilor folosite. Pentru fiecare decizie sunt prezentate și criteriile care au stat la baza ei, cu scopul de a facilita înțelegerea și dezvoltarea ulterioară a aplicației.

\section{Arhitectura generală a soluției}\label{sec:arhitectura}

\subsection{Principii de organizare}\label{subsec:principii}

Organizarea codului respectă trei principii, impuse de natura experimentală a lucrării.

\textbf{Separarea logicii de execuție.} Întreaga logică se află într-un pachet de biblioteci, iar scripturile executabile din rădăcina proiectului conțin doar apelul etapei corespunzătoare. Un script are, în medie, sub o sută de linii și nu implementează nimic propriu. Consecința practică este că orice funcționalitate poate fi apelată independent, atât din alt script, cât și din interfața interactivă, fără duplicarea codului.

\textbf{Etape reluabile independent.} Fiecare etapă își persistă rezultatele pe disc într-un format tabelar, iar etapa următoare le citește de acolo. O etapă poate fi astfel reluată fără a le relua pe cele anterioare, ceea ce este esențial într-un context în care antrenarea modelelor durează zeci de minute, iar analizele ulterioare sunt reluate frecvent.

\textbf{Model unic de acces la predicție.} Toate predicțiile trec printr-un singur modul, indiferent de model și de etapă. Motivul este detaliat în subsecțiunea~\ref{subsec:inferenta}: fără această centralizare, două etape ar putea obține predicții diferite pentru același flux, ceea ce ar invalida măsurătoarea.

\subsection{Structura pe pachete}\label{subsec:pachete}

Codul este organizat în pachetele prezentate în tabelul~\ref{tab:pachete}. Gruparea urmează etapele conceptuale din secțiunea~\ref{sec:structura}, nu tipurile de obiecte, astfel încât modificarea unei etape să fie localizată într-un singur director.

\begin{table}[H]
    \caption{Organizarea codului pe pachete și responsabilitățile acestora}
    \centering
    \begin{tabular}{|l|p{8.4cm}|}
        \hline
        \textbf{Pachet} & \textbf{Responsabilitate} \\
        \hline
        Nucleu & Configurație, încărcarea datelor, preprocesare, definirea modelelor, evaluare, analiză exploratorie, utilitare comune \\
        \hline
        Inferență & Punct unic de predicție pentru toate modelele \\
        \hline
        Rețea neuronală & Arhitectura, preprocesarea specifică, bucla de antrenare, adaptorul de predicție \\
        \hline
        Perturbare & Schema caracteristicilor, generatoarele de transformări, propagarea dependențelor, componenta de validare \\
        \hline
        Evaziune & Măsurarea rezultatelor, taxonomia, agregarea, generarea figurilor \\
        \hline
        Sensibilitate & Importanța prin permutare, corelarea cu evaziunea, analiza destinațiilor \\
        \hline
        Antrenare adversarială & Construirea seturilor augmentate, grupurile de antrenare,
        evaluarea pe mulțimea comună, fișa de reproductibilitate \\
        \hline
        Interfață & Aplicația interactivă și adaptorul de perturbare live \\
        \hline
    \end{tabular}
    \label{tab:pachete}
\end{table}

Structura de dependențe este reprezentată în figura~\ref{fig:pachete}.

Dependențele sunt, cu o singură excepție, unidirecționale: pachetele de nivel superior importă din cele inferioare, niciodată invers. Scripturile de rulare nu sunt importate de nimeni, iar pachetele de interfață și de antrenare adversarială sunt și ele consumatori finali. Antrenarea adversarială este singura care produce modele noi, motiv pentru care scrie exclusiv în directoare proprii.

Excepția merită enunțată, pentru că o diagramă care ar ascunde-o ar descrie altceva decât codul. Modulul de inferență, aflat în nucleu, are nevoie de clasa care încarcă rețeaua neuronală, iar aceasta din urmă importă la rândul ei din nucleu, deci cele două pachete se referă reciproc. Ciclul nu produce însă o eroare la încărcare, întrucât importul din nucleu către rețea se execută în interiorul funcției care îl folosește, nu la nivelul modulului. Alegerea are și un al doilea motiv, practic: etapele care lucrează exclusiv cu ansamblurile de arbori nu încarcă astfel biblioteca de învățare profundă, care este de departe cea mai costisitoare dependență a proiectului.

\begin{figure}[!h]
    \centering
    \begin{tikzpicture}[node distance=8mm and 9mm]
        \node[bloc, minimum width=112mm] (scripturi)
            {\textbf{Scripturi de rulare} \;---\; câte unul pentru fiecare etapă};

        \node[bloc, minimum width=26mm, below=10mm of scripturi] (adv)
            {\textbf{Antrenare}\\ \textbf{adversarială}};
        \node[bloc, minimum width=26mm, left=18mm of adv] (ui) {\textbf{Interfață}};
        \node[bloc, minimum width=26mm, right=18mm of adv] (sens) {\textbf{Sensibilitate}};

        \node[bloc, minimum width=26mm, below=10mm of adv] (evaz) {\textbf{Evaziune}};
        \node[bloc, minimum width=26mm, below=10mm of ui] (pert) {\textbf{Perturbare}};
        \node[bloc, minimum width=26mm, below=10mm of sens] (retea)
            {\textbf{Rețea}\\ \textbf{neuronală}};

        \node[bloc, minimum width=112mm, below=10mm of evaz] (nucleu)
            {\textbf{Nucleu} \;-\; configurație, date, preprocesare, modele, evaluare,\\
             utilitare și \textbf{inferența} (punctul unic de predicție)};

        \draw[sageata] (scripturi.south -| ui) -- (ui);
        \draw[sageata] (scripturi) -- (adv);
        \draw[sageata] (scripturi.south -| sens) -- (sens);

        \draw[sageata] (ui) -- (pert);
        \draw[sageata] (ui.south east) -- (evaz.north west);
        \draw[sageata] (adv) -- (evaz);
        \draw[sageata] (adv.south west) -- (pert.north east);
        \draw[sageata] (adv.south east) -- (retea.north west);
        \draw[sageata] (sens.south west) -- (evaz.north east);
        \draw[sageata] (retea) -- (nucleu.north -| retea);
        \draw[sageata] (evaz) -- (nucleu.north -| evaz);
        \draw[sageata] (pert) -- (nucleu.north -| pert);
        \draw[sageata] (evaz) -- (pert);
        \draw[sageata] (retea.south west) to[out=200, in=340] (pert.south east);

        \draw[punctat] (nucleu.east) -- ++(9mm,0) |-
            node[eticheta, pos=0.25, right] {import întârziat} (retea.east);
    \end{tikzpicture}
    \caption{Dependențele reale dintre pachete, extrase din instrucțiunile de import. Săgețile continue merg de sus în jos: fiecare pachet importă din cele situate sub el. Singura excepție este săgeata punctată, prin care nucleul ajunge la rețeaua neuronală; ea închide un ciclu, evitat în practică prin efectuarea importului abia în momentul apelului, nu la încărcarea modulului.}
    \label{fig:pachete}
\end{figure}

\subsection{Configurația centralizată}\label{subsec:configuratie}

Toate constantele experimentului sunt declarate într-un singur modul de configu\-rație: căile către date și rezultate, \textit{seed}-ul generatorului de numere aleatoare, hiperparametrii celor trei modele, taxonomia caracteristicilor din tabelul~\ref{tab:taxonomie}, nivelurile de perturbare, pragurile de validitate fizică și denumirile categoriilor de rezultat.

Într-un experiment în care aceeași constantă este folosită de mai multe etape rulate la momente diferite, o valoare duplicată în două locuri conduce la rezultate care par valide, dar provin din configurații diferite. Un exemplu concret îl constituie valoarea TTL asociată traficului legitim, folosită simultan de generatorul de perturbări, de mecanismul de rezolvare a contorului de context și de interfața interactivă. Declararea ei într-un singur loc garantează că cele trei componente se referă la aceeași valoare.

Modulul conține, în plus, comentarii extinse care documentează motivul fiecărei alegeri neevidente, de exemplu, de ce nivelurile TTL sunt exact cele alese, sau de ce predicția se execută pe un singur fir de execuție. Aceste comentarii constituie, în practică, jurnalul deciziilor experimentale.

\subsection{Fluxul de execuție}\label{subsec:flux}

Etapele se execută în ordinea din tabelul~\ref{tab:etape}. Ordinea este impusă de dependențele de date: fiecare etapă consumă ieșirile celei precedente.

\begin{table}[H]
    \caption{Etapele de execuție, dependențele și rezultatele produse}
    \centering
    \begin{tabular}{|l|p{3.3cm}|p{4.2cm}|}
        \hline
        \textbf{Etapă} & \textbf{Consumă} & \textbf{Produce} \\
        \hline
        Antrenarea ansamblurilor & Setul de date brut & Modele salvate, metrici, predicții de referință \\
        \hline
        Antrenarea rețelei & Setul de date brut & Model salvat, metrici, predicții de referință \\
        \hline
        Generarea variantelor & Predicțiile de referință & Manifest, raport de validitate, eșantioane \\
        \hline
        Măsurarea evaziunii & Modelele, variantele & Rate, intervale, rezultate per flux \\
        \hline
        Analiza sensibilității & Modelele, rezultatele per flux & Importanțe, corelații, figuri \\
        \hline
        Antrenarea adversarială & Setul de date brut, generatoarele de variante & Modele per grup, mulțimea comună, verdictul față de criterii \\
        \hline
        Interfața & Toate cele anterioare & Explorare interactivă \\
        \hline
    \end{tabular}
    \label{tab:etape}
\end{table}

Etapa de antrenare a rețelei neuronale este independentă de cea a ansamblurilor de arbori și poate fi executată în orice ordine față de aceasta. Ambele trebuie însă să preceadă generarea variantelor, deoarece mulțimea eligibilă se calculează din predicțiile de referință ale tuturor modelelor.

Antrenarea adversarială ocupă o poziție aparte în această succesiune. Ea consumă datele brute și mecanismul de generare, nu rezultatele etapelor de măsurare, deci ar putea fi executată imediat după prima etapă. Este plasată totuși la sfârșit, din două motive: comparația pe care o produce are sens numai raportată la vulnerabilitatea deja măsurată, iar criteriile după care se judecă rezultatul se formulează pornind de la valorile obținute în etapele anterioare. Fiind singura etapă care produce modele noi, ea scrie exclusiv în directoare proprii, astfel încât toate rezultatele precedente rămân valabile după rulare.

\section{Încărcarea și pregătirea datelor}\label{sec:date}

\subsection{Setul de date și partiționarea}\label{subsec:setdate}

Se utilizează setul UNSW-NB15 \cite{moustafa2015unsw}, în partiția oficială pusă la dispoziție de autori, cu 175\,341 de înregistrări pentru antrenare și 82\,332 pentru testare. Partiția oficială este preferată unei împărțiri proprii din două motive: rezultatele devin comparabile cu cele raportate în literatură, iar împărțirea aleatoare a unui set de trafic poate plasa fluxuri provenite din aceeași sesiune în ambele partiții, producând o scurgere de informație greu de detectat.

Fiecare înregistrare conține 45 de coloane: un identificator, 42 de caracteristici predictive (39 numerice și 3 categorice) și două etichete, una binară și una cu categoria de atac. Cifra de 42 folosită în toată lucrarea se referă la caracteristicile predictive, adică la coloanele rămase după excluderea identificatorului și a celor două etichete. Se folosește exclusiv eticheta multi-clasă, cu zece valori posibile. Coloana de identificator este păstrată separat de caracteristici: ea nu intră niciodată în model, dar este necesară pentru a putea urmări un flux individual de-a lungul întregului lanț de prelucrare, de la predicția de referință până la rezultatul perturbării.

Modulul de încărcare validează la fiecare apel prezența coloanelor așteptate și consistența tipurilor, semnalând o eroare explicită în locul unei erori tardive apărute în etapa de antrenare.

\subsection{Analiza exploratorie}\label{subsec:eda}

Etapa de analiză exploratorie produce caracterizările numerice și grafice ale setului de date. Ea raportează distribuția claselor și raportul de dezechilibru, cardinalitatea variabilelor categorice împreună cu numărul de categorii aflate sub pragul de raritate, tabelul asimetriei distribuțiilor numerice și corelațiile dintre caracteristicile numerice și etichetă.

Trei figuri sunt generate și salvate pentru a fi citate în capitolul de rezultate: distribuția claselor pe scară logaritmică, corelațiile cu eticheta și histogramele caracteristicilor cu cea mai pronunțată asimetrie. Scara logaritmică este necesară pentru prima figură, întrucât raportul dintre clasa cea mai frecventă și cea mai rară face ca, pe scară liniară, clasele minoritare să fie invizibile.

Transformarea logaritmică aplicată histogramelor este folosită \emph{exclusiv} pentru lizibilitatea reprezentării grafice și nu intră în fluxul de modelare. Motivul este cel discutat în subsecțiunea~\ref{subsec:arbori}: o transformare monotonă aplicată independent fiecărei caracteristici nu modifică deciziile unui ansamblu de arbori, deci nu ar avea niciun efect asupra rezultatelor, iar aplicarea ei doar pentru unele modele ar introduce o diferență inutilă între condițiile experimentale.

\subsection{Gruparea categoriilor rare}\label{subsec:rare}

Variabila care descrie protocolul are 133 de valori distincte, cu o distribuție puternic dezechilibrată: primele două valori acoperă împreună peste jumătate din observații, în timp ce 118 dintre ele apar de mai puțin de 200 de ori fiecare. Codificarea directă prin variabile indicator produce astfel o matrice cu peste o sută de coloane aproape constante, ceea ce ridică riscul supraînvățării unor categorii cu suport statistic redus.

S-a implementat, deci, un transformator care grupează categoriile rare. La ajustare, el reține pentru fiecare coloană categorică mulțimea valorilor care apar de cel puțin un număr prestabilit de ori în datele de antrenare. La transformare, orice valoare din afara acestei mulțimi este înlocuită cu o etichetă rezervată. Pragul folosit este de 50 de apariții. Transformatorul este integrat în fluxul de preprocesare \emph{înaintea} codificării, astfel încât să fie salvat împreună cu modelul și să se aplice identic la reîncărcare.

\begin{verbatim}
class RareCategoryGrouper(BaseEstimator, TransformerMixin):
    """Inlocuieste categoriile rare sau necunoscute."""

    def fit(self, X, y=None):
        # Multimea retinuta se calculeaza EXCLUSIV pe datele primite
        # la ajustare, adica pe partitia de antrenare.
        self.keep_ = {
            col: set(X[col].value_counts()
                     [lambda s: s >= self.min_freq].index)
            for col in self.columns
        }
        return self

    def transform(self, X):
        X = X.copy()
        for col in self.columns:
            # Acelasi test prinde si categoriile rare din train,
            # si pe cele nevazute: ambele cad in afara multimii.
            X[col] = X[col].where(X[col].isin(self.keep_[col]),
                                  self.other_label)
        return X
\end{verbatim}

Separarea dintre cele două metode este esențială. Mulțimea valorilor păstrate
se stabilește în \texttt{fit}, exclusiv din datele de antrenare, și rămâne
fixată după aceea. Partiția de testare trece doar prin \texttt{transform}, deci
nu poate influența ce categorii sunt considerate frecvente. O implementare care
ar recalcula mulțimea la fiecare transformare ar produce o scurgere de
informație dinspre setul de testare, greu de observat în rezultate.

Al doilea aspect este că nu există o ramură separată pentru categoriile
nevăzute. Un singur test, apartenența la \texttt{keep\_}, tratează deopotrivă o
categorie rară din antrenare și o valoare care apare pentru prima dată la
inferență, întrucât niciuna dintre ele nu se află în mulțimea învățată.

Efectul cantitativ al grupării trebuie raportat corect, întrucât este mai mic decât sugerează cardinalitatea inițială. La pragul ales, doar trei protocoale sunt colapsate, iar lățimea codificării acestei coloane scade de la 133 la 131 de poziții. Motivul este că lungul șir de categorii rare nu coboară, în general, sub pragul de 50: majoritatea protocoalelor apar de ordinul a o sută de ori, ceea ce reflectă modul sintetic în care a fost generat traficul. Reducerea dimensionalității nu constituie, prin urmare, principalul beneficiu al mecanismului în acest set de date.

Variabila care descrie starea conexiunii conține în partiția de testare două valori care nu apar deloc în cea de antrenare. Fără o tratare explicită, codificatorul le-ar transforma tăcut într-un vector nul, comportament corect din punct de vedere tehnic, dar complet invizibil în raportare. Prin grupare, aceste valori sunt încadrate explicit în categoria rezervată, iar situația devine observabilă și documentată. Mecanismul a fost păstrat, prin urmare, pentru corectitudinea și transparența tratării categoriilor nevăzute, nu pentru compresia spațiului de caracteristici.

\subsection{Fluxul de preprocesare}\label{subsec:preprocesare}

Preprocesarea este încapsulată într-un flux compus, ceea ce garantează că aceleași transformări se aplică identic la antrenare și la predicție. Structura este comună celor două ansambluri de arbori: gruparea categoriilor rare, apoi codificarea prin variabile indicator a coloanelor categorice, cu transmiterea nemodificată a celor numerice.

Absența normalizării valorilor numerice este o decizie deliberată, justificată în subsecțiunea~\ref{subsec:arbori}. Codificatorul păstrează, suplimentar, tratarea valorilor nevăzute ca plasă de siguranță, chiar dacă gruparea prealabilă ar trebui să le elimine.

Rețeaua neuronală folosește un flux diferit, descris în subsecțiunea~\ref{subsec:retea}. Diferența privește exclusiv reprezentarea internă a valorilor, în timp ce mulțimea caracteristicilor, împărțirea în partiții și etichetele rămân identice.

\section{Implementarea clasificatoarelor}\label{sec:clasificatoare}

\subsection{Ansamblurile de arbori}\label{subsec:ansambluri}

Cele două ansambluri sunt construite peste fluxul de preprocesare comun. Random Forest folosește 200 de arbori și ponderarea claselor invers proporțională cu frecvența lor. Modelul cu XGBoost folosește 300 de arbori, adâncime maximă 8, rată de învățare 0{,}1 și eșantionare a rândurilor la fiecare iterație.

Ponderarea claselor necesită tratamente diferite la cele două biblioteci. În primul caz este disponibil un parametru dedicat. În al doilea, ponderile se calculează explicit și se transmit ca vector de ponderi per observație la apelul de antrenare. Rezultatul este echivalent, dar diferența de interfață trebuie cunoscută la întreținere.

O particularitate de implementare merită semnalată: biblioteca de XGBoost necesită etichete codificate numeric, în timp ce restul lanțului lucrează cu etichete textuale. Conversia este încapsulată împreună cu modelul, iar maparea inversă este salvată alături de acesta, astfel încât predicțiile expuse în exterior să fie întotdeauna textuale. Fără această încapsulare, ordinea claselor ar deveni o convenție implicită, ușor de încălcat la o modificare ulterioară.

După antrenare, fiecare model este salvat pe disc, apoi reîncărcat și verificat: predicțiile modelului reîncărcat trebuie să coincidă cu cele ale modelului din memorie. Verificarea prinde erorile de serializare, care altfel s-ar manifesta abia în etapele următoare, sub forma unor rezultate inexplicabile.

\subsection{Transformer}\label{subsec:retea}

Mecanismul de atenție este scris de la zero, nu preluat din componentele gata făcute ale bibliotecii de învățare profundă. Alegerea costă mai mult cod, dar aduce două lucruri pe care varianta scurtă nu le oferă: fiecare operație descrisă în subsecțiunea~\ref{subsec:atentie} apare explicit și poate fi verificată separat, iar ponderile de atenție rămân accesibile din exterior, în loc să rămână închise într-o componentă opacă. A doua proprietate este cea care face posibilă inspectarea a ceea ce modelul consideră relevant într-un flux dat.

\textbf{Tokenizatorul de caracteristici} transformă vectorul de intrare într-o secvență de reprezentări. Fiecare coloană numerică primește un vector de ponderi și unul de deplasare proprii, astfel încât valoarea scalară a coloanei este proiectată într-un spațiu de dimensiune 192 printr-o transformare afină învățată separat pentru fiecare coloană. Fiecare coloană categorică primește un tabel de reprezentări indexat după valoarea categoriei. Se adaugă un vector de agregare, inițializat aleator și învățat, care va colecta informația din întreaga secvență. Inițializarea tuturor parametrilor este uniformă, cu limite invers proporționale cu rădăcina dimensiunii reprezentării, ca în lucrarea de referință \cite{gorishniy2021revisiting}.

Structura completă a rețelei este prezentată în figura următoare, iar implementarea mecanismului de atenție este reprodusă în secțiunea~\ref{sec:anexa-atentie}.

\begin{figure}[!ht]
    \centering
    \begin{tikzpicture}[node distance=5mm and 6mm]
        \node[bloc, minimum width=26mm] (num) {39 coloane\\ numerice};
        \node[bloc, minimum width=26mm, right=10mm of num] (cat) {3 coloane\\ categorice};
        \node[bloc, minimum width=26mm, above=of num] (proj) {Proiecție afină\\ per coloană};
        \node[bloc, minimum width=26mm, above=of cat] (emb) {Tabel de\\ reprezentări};
        \node[bloc, minimum width=16mm, left=10mm of proj] (cls) {token\\ \texttt{[CLS]}};
        \node[bloc, minimum width=74mm, above=6mm of proj.north east, anchor=south] (seq)
            {Secvență de $1 + 42$ tokeni, fiecare de dimensiune $d = 192$};
        \node[bloc, minimum width=74mm, above=of seq] (blocuri)
            {$3 \times$ \; \{ atenție multi-head (8 capete) $\rightarrow$ rețea complet conectată \}\\
             fiecare cu normalizare și conexiune reziduală};
        \node[bloc, minimum width=40mm, above=of blocuri] (head) {Strat liniar aplicat\\ reprezentării \texttt{[CLS]}};
        \node[bloc, minimum width=40mm, above=of head] (out) {10 clase};

        \draw[sageata] (num) -- (proj);
        \draw[sageata] (cat) -- (emb);
        \draw[sageata] (proj.north) -- ++(0,3mm) -| ($(seq.south)+(-10mm,0)$);
        \draw[sageata] (emb.north) -- ++(0,3mm) -| ($(seq.south)+(10mm,0)$);
        \draw[sageata] (cls.north) |- ($(seq.west)+(-2mm,0)$) -- (seq.west);
        \draw[sageata] (seq) -- (blocuri);
        \draw[sageata] (blocuri) -- (head);
        \draw[sageata] (head) -- (out);
    \end{tikzpicture}
    \caption{Arhitectura modelului Transformer. Fiecare caracteristică devine un token, iar clasificarea se face din reprezentarea finală a tokenului de agregare.}
    \label{fig:transformer}
\end{figure}

\textbf{Blocul de transformare} aplică succesiv atenție multi-head cu 8 capete și o rețea complet conectată, fiecare precedată de normalizare și urmată de o conexiune reziduală. Se folosesc trei astfel de blocuri. Regularizarea se aplică diferențiat: 0{,}2 pe ponderile de atenție și 0{,}1 în rețeaua complet conectată, valorile fiind alese astfel încât capacitatea modelului să nu fie limitată excesiv pe clasele rare. Clasificarea se realizează dintr-un strat liniar aplicat reprezentării finale a vectorului de agregare. Modelul rezultat are aproximativ 789 de mii de parametri.

\textbf{Preprocesarea specifică} diferă de cea a ansamblurilor prin două aspecte. Valorile numerice sunt standardizate, din motivele expuse în subsecțiunea~\ref{subsec:atentie}. Standardizarea este ajustată exclusiv pe subsetul efectiv de antrenare, cu excluderea feliei de validare, pentru ca aceasta din urmă să rămână neinfluențată și criteriul de oprire timpurie să fie nepărtinitor. Coloanele categorice sunt convertite în indici întregi, nu în variabile indicator, întrucât sunt consumate de tabelele de reprezentări.

\textbf{Bucla de antrenare} folosește optimizatorul AdamW cu rată de învățare $10^{-4}$, regularizare $10^{-5}$ și loturi de 512 observații, pe cel mult 100 de epoci. Funcția de cost este entropia încrucișată cu ponderi per clasă, calculate din frecvențele din setul de antrenare. Se rezervă 15\% din datele de antrenare, stratificat, pentru validare.

\textbf{Selecția modelului final} constituie punctul care a necesitat cea mai atentă proiectare. Criteriul natural, păstrarea modelului cu cel mai bun scor de validare, s-a dovedit nepotrivit în acest context. Scorul macro-$F1$ de validare oscilează puternic de la o epocă la alta, deoarece este dominat de clasele rare, unde câteva observații clasificate diferit modifică sensibil media neponderată din secțiunea~\ref{sec:metrici}. Selecția pe valoarea instantanee riscă astfel să rețină un maxim conjunctural, apărut din fluctuație, în locul unui model aflat pe o traiectorie stabil ascendentă.

Soluția adoptată este selecția pe scorul \emph{netezit}: la fiecare epocă se calculează media scorurilor de validare pe o fereastră mobilă de trei epoci, iar modelul este reținut atunci când această medie se îmbunătățește. Același criteriu netezit guvernează și oprirea timpurie, cu o răbdare de 10 epoci. Procedura nu garantează un scor final mai bun, o fluctuație favorabilă poate produce, punctual, o valoare superioară, dar garantează că modelul reținut este reprezentativ pentru o regiune stabilă a traiectoriei de antrenare, ceea ce este cerința corectă atunci când modelul urmează să fie supus unor experimente ulterioare.

Antrenarea este configurată să folosească algoritmi determiniști, cu prețul unui timp de execuție cu aproximativ 7\% mai mare. Fără această constrângere, operațiile de propagare inversă prin tabelele de reprezentări produc rezultate ușor diferite între rulări, iar modelul obținut nu ar fi reproductibil.

\subsection{Stratul de inferență deterministă}\label{subsec:inferenta}

Toate predicțiile din lucrare trec printr-un modul unic, care expune aceeași interfață pentru cele trei modele și ascunde diferențele dintre biblioteci. Existența lui rezolvă o problemă descoperită experimental și care, nesoluționată, ar fi compromis măsurătoarea.

Random Forest produce, pentru un număr mic de fluxuri, predicții care variază între apeluri identice. Cauza este că decizia se ia prin însumarea voturilor arborilor, iar când primele două clase primesc un număr foarte apropiat de voturi, ordinea de însumare, care depinde de repartizarea pe fire de execuție, determină care dintre ele are valoarea finală mai mare. În setul analizat există 44 de fluxuri aflate în această situație, dintre care 43 la egalitate exactă și unul cu o diferență la limita preciziei reprezentării în virgulă mobilă.

Consecința asupra experimentului este disproporționată față de numărul de fluxuri afectate. Dacă predicția de referință și predicția pe varianta perturbată se obțin din apeluri diferite, un astfel de flux poate fi raportat ca schimbându-și clasa fără ca nimic să fi fost modificat, iar schimbarea ar fi atribuită eronat perturbării. Într-un context în care unele clase de atac au sub o sută de observații, câteva astfel de comutări false nu se pierd în zgomot.

Soluția este forțarea execuției pe un singur fir la încărcarea modelului, ceea ce elimină complet variația. Costul este o creștere a timpului de predicție de la aproximativ 0{,}4 la 1{,}2 secunde pentru întregul set de atac, neglijabilă la scara experimentului. Este important de subliniat că măsura nu modifică niciun rezultat: predicțiile obținute pe un singur fir coincid cu cele stocate anterior. Ea elimină variabilitatea, nu introduce o valoare diferită.

Modulul realizează și adaptarea rețelei neuronale la aceeași interfață, aplicând intern preprocesarea specifică și convertind ieșirile numerice în etichete textuale. Etapele ulterioare pot astfel trata cele trei modele uniform, fără ramificații condiționale.

\section{Motorul de generare a variantelor perturbate}\label{sec:motor}

\subsection{Schema caracteristicilor}\label{subsec:schema}

Un modul dedicat menține corespondența dintre numele coloanelor și categoriile din tabelul~\ref{tab:taxonomie} și expune funcții de interogare folosite de generatoare și de validatoare. Centralizarea garantează că generatorul și validatorul operează cu aceeași definiție a mulțimii caracteristicilor imuabile. În absența ei, o modificare a taxonomiei într-un singur loc ar produce un validator care aprobă variante pe care ar trebui să le respingă.

Modulul asigură și restaurarea tipurilor de date. Operațiile de propagare produc valori în virgulă mobilă chiar și pentru coloane care în setul original sunt întregi, iar diferența de tip poate modifica comportamentul codificatorului. După fiecare transformare, tipurile sunt readuse la cele ale cadrului original.

\subsection{Generatoarele de transformări}\label{subsec:generatoare}

Fiecare tip de transformare este implementat ca o funcție cu semnătură comună, care primește cadrul de intrare, nivelul de intensitate și contextul precalculat, și returnează cadrul perturbat împreună cu metadatele transformării. Uniformitatea semnăturii permite tratarea generatoarelor ca intrări într-un registru, iar adăugarea unui nou tip de transformare nu necesită modificarea codului care le apelează. Funcția care orchestrează pașii este reprodusă în secțiunea~\ref{sec:anexa-generare}.

\textbf{Transformarea TTL} atribuie valoarea TTL a pachetelor emise o valoare-țintă și declanșează rezolvarea contorului de context conform politicii selectate. Valorile-țintă sunt 64, 62 și 31, alese dintre valorile efectiv observate în captură. O valoare precum 128, deși uzuală în alte contexte, nu apare deloc în acest set și ar produce o combinație pe care modelul nu a întâlnit-o niciodată, ceea ce ar confunda efectul perturbării cu efectul extrapolării.

\textbf{Adăugarea de umplutură} (\verb|padding|) crește volumul de octeți transmiși cu o fracțiune dată. Implementarea tratează explicit limita fizică a dimensiunii pachetului: dacă mărirea pachetelor existente ar depăși unitatea maximă de transmisie, generatorul adaugă pachete suplimentare în loc să depășească limita, ceea ce modifică în mod corect și numărul de pachete, deci și mărimile derivate din el. Plafonul per pachet se calculează relativ la valoarea originală, deoarece setul conține o înregistrare a cărei dimensiune medie o depășește deja pe cea nominală.

\textbf{Dilatarea temporală} (\verb|timing|) multiplică durata fluxului cu un factor supraunitar, ceea ce antrenează prin propagare atât mărimile temporale ale sursei, cât și pe cele ale destinației, conform discuției din subsecțiunea~\ref{subsec:taxonomie}.

\textbf{Reducerea contoarelor de conexiuni} (\verb|connection_rate|) le scalează descrescător, cu un prag inferior egal cu valoarea minimă observată în date. Nivelul cel mai intens coboară contoarele direct la acest minim.

\textbf{Transformările combinate} aplică succesiv modificarea TTL, umplutura și dilatarea temporală, la intensități crescătoare, propagarea fiind executată o singură dată, la final, pentru a evita compunerea repetată a rapoartelor.

\textbf{Varianta de referință cu constrângeri relaxate} modifică suplimentar valoarea TTL a pachetelor de răspuns. Ea este marcată explicit ca nerealizabilă atât în cod, cât și în fișierele de rezultate, și este exclusă din agregările principale. Validatorul primește, pentru această singură variantă, o excepție declarată explicit, astfel încât modificarea unei coloane altfel imuabile să fie permisă doar aici și niciodată tacit.

Setul complet de variante generate este prezentat în tabelul~\ref{tab:grila}. Fiecare linie corespunde unui tip de transformare, iar numărul de niveluri determină câte variante distincte se produc.

\begin{table}[!ht]
    \caption{Grila de variante generate}
    \centering
    \begin{tabular}{|l|c|p{5.4cm}|}
        \hline
        \textbf{Tip} & \textbf{Niveluri} & \textbf{Parametri} \\
        \hline
        Identitate & 1 & Nicio modificare (control) \\
        \hline
        TTL, politică conservatoare & 3 & Valori-țintă 64, 62, 31 \\
        \hline
        TTL, politică optimistă & 3 & Valori-țintă 64, 62, 31 \\
        \hline
        Umplutură & 4 & $+10\%$, $+25\%$, $+50\%$, $+100\%$ \\
        \hline
        Temporizare & 4 & $\times 1{,}5$, $\times 2$, $\times 5$, $\times 10$ \\
        \hline
        Rată de conexiuni & 4 & $\times 0{,}75$, $\times 0{,}5$, $\times 0{,}25$, minim \\
        \hline
        Combinată, conservatoare & 4 & Intensități crescătoare \\
        \hline
        Combinată, optimistă & 4 & Intensități crescătoare \\
        \hline
        Referință relaxată & 1 & \textbf{Nerealizabilă}: atinge și TTL-ul victimei \\
        \hline
        \multicolumn{2}{|l|}{\textbf{Total}} & \textbf{28 de variante} \\
        \hline
    \end{tabular}
    \label{tab:grila}
\end{table}

Nivelul 0 este întotdeauna varianta identică și servește drept control, conform invariantului identității din subsecțiunea~\ref{subsec:rezultate}. Cele 26 de variante realizabile de la nivelurile nenule constituie baza analizei de corelație din secțiunea~\ref{sec:implsensibilitate}.

O transformare a fost eliminată în cursul dezvoltării. Modificarea parametrilor ferestrei TCP s-a dovedit lipsită de conținut: analiza datelor a arătat că această caracteristică ia o singură valoare pentru toate fluxurile care folosesc protocolul TCP și o altă valoare unică pentru restul, fără excepție. Ea nu codifică, prin urmare, nimic în plus față de protocolul însuși, iar perturbarea ei fie nu ar schimba nimic, fie ar contrazice protocolul, care este o caracteristică imuabilă. Eliminarea unei transformări lipsite de conținut este preferabilă raportării unei rate de evaziune nule obținute dintr-un motiv greșit.

\subsection{Propagarea dependențelor}\label{subsec:implpropagare}

Propagarea implementează direct observația centrală din subsecțiunea~\ref{subsec:propagare}. După ce generatorul a stabilit valorile de bază, modulul de dependențe calculează rapoartele dintre valorile noi și cele originale și le aplică multiplicativ caracteristicilor derivate, conform relațiilor din tabelul~\ref{tab:propagare}.

Tratarea numitorilor nuli este esențială pentru robustețea implementării. Împărți\-rea se execută cu o valoare de rezervă egală cu unitatea acolo unde numitorul este zero sau rezultatul nu este finit, ceea ce lasă caracteristica derivată neschimbată. Modulul raportează, pentru fiecare variantă, numărul de rânduri afectate de fiecare tip de caz degenerat, fluxuri fără octeți transmiși, cu durată nulă sau cu un singur pachet, astfel încât amploarea neutralizării să fie vizibilă în fișa rulării, nu ascunsă.

O atenție deosebită necesită intervalul mediu dintre pachete, a cărui propagare depinde de numărul de \emph{intervale}, adică de numărul de pachete minus unu. Pentru un flux cu un singur pachet, această mărime este nulă, iar raportul devine nedefinit, iar cazul este tratat separat, nu prin formula generală. Codul corespunzător se află în secțiunea~\ref{sec:anexa-propagare}.

\subsection{Rezolvarea contorului de context}\label{subsec:ct}

Contorul care combină starea conexiunii cu valorile TTL nu poate fi recalculat, conform argumentului din subsecțiunea~\ref{subsec:interval}. Implementarea materializează cele două politici printr-un mecanism cu trei niveluri de rezervă.

Se construiește mai întâi, din datele disponibile, o corespondență între tripletul format din starea conexiunii și cele două valori TTL, pe de o parte, și valoarea dominantă a contorului, pe de altă parte. Politica conservatoare păstrează valoarea originală și nu consultă această corespondență. Politica optimistă caută tripletul rezultat în urma perturbării, iar dacă acesta a fost observat, preia valoarea asociată. Dacă tripletul nu apare niciodată în date, se consultă o a doua corespondență, construită exclusiv din traficul legitim, care asociază fiecărei valori TTL semnătura tipică a traficului normal. Dacă nici aceasta nu conține valoarea căutată, se recurge la o constantă declarată explicit.

Este important de precizat că această corespondență se construiește din reuniunea partițiilor de antrenare și testare. Alegerea nu constituie o scurgere de informație către model, întrucât nu se antrenează nimic pe ea: corespondența modelează comportamentul \emph{senzorului} care a produs datele, adică o proprietate a procesului de generare, nu o proprietate a modelului evaluat. Restrângerea ei la partiția de antrenare ar produce o modelare mai slabă a aceluiași proces, fără niciun câștig metodologic.

Manifestul înregistrează, pentru fiecare variantă, câte rânduri au fost rezolvate prin fiecare dintre cele trei niveluri, ceea ce permite evaluarea gradului în care rezultatul depinde de nivelurile de rezervă. Mecanismul cu trei niveluri este reprodus în secțiunea~\ref{sec:anexa-ct}.

\subsection{Componenta de validare}\label{subsec:validare}

Fiecare variantă generată este supusă unui set de verificări înainte de a fi utilizată. Verificările sunt prezentate în tabelul~\ref{tab:verificari}.

\begin{table}[H]
    \caption{Verificările aplicate fiecărei variante perturbate}
    \centering
    \begin{tabular}{|l|p{7.6cm}|}
        \hline
        \textbf{Verificare} & \textbf{Conținut} \\
        \hline
        Direcționalitate & Octeții și durata pot doar crește, contoarele pot doar scădea \\
        \hline
        Limite de domeniu & Valorile rămân în intervalele fizic realizabile \\
        \hline
        Imuabilitate & Coloanele victimei și cele identitare rămân neschimbate \\
        \hline
        Consistența dependențelor & Fiecare caracteristică derivată rămâne în acord cu mărimile din care se calculează, flux cu flux \\
        \hline
        Finitudine & Propagarea nu a introdus valori nedefinite sau infinite \\
        \hline
        Schemă & Coloanele, ordinea și tipurile corespund celor așteptate \\
        \hline
    \end{tabular}
    \label{tab:verificari}
\end{table}

Poziția validării în lanțul de generare este prezentată în figura~\ref{fig:motor}. Verificarea este ultima etapă înaintea utilizării, iar o variantă respinsă nu este corectată tacit, ci oprește execuția.

\begin{figure}[!ht]
    \centering
    \begin{tikzpicture}[node distance=6mm and 18mm]
        \node[bloc, minimum width=24mm] (in) {Fluxuri\\ eligibile};
        \node[bloc, minimum width=26mm, right=of in] (gen) {Generator\\ (tip, nivel)};
        \node[bloc, minimum width=26mm, right=of gen] (prop) {Propagarea\\ dependențelor};
        \node[bloc, minimum width=26mm, below=11mm of prop] (ct) {Rezolvarea\\ contorului};
        \node[bloc, minimum width=26mm, left=of ct] (tip) {Restaurarea\\ tipurilor};
        \node[bloc, minimum width=26mm, left=of tip] (val) {Validare\\ (6 verificări)};
        \node[bloc, minimum width=30mm, below=12mm of val, xshift=-6mm] (stop) {\textit{încălcare:}\\ execuție oprită};
        \node[bloc, minimum width=30mm, below=12mm of tip, xshift=6mm] (ok) {\textit{conformă:}\\ variantă utilizabilă};

        \draw[sageata] (in) -- (gen);
        \draw[sageata] (gen) -- (prop);
        \draw[sageata] (prop) -- (ct);
        \draw[sageata] (ct) -- (tip);
        \draw[sageata] (tip) -- (val);
        \draw[sageata] (val.south) -- (stop.north);
        \draw[sageata] (val.south east) -- (ok.north west);
    \end{tikzpicture}
    \caption{Fluxul motorului de generare a variantelor perturbate. Validarea este ultima etapă înaintea utilizării: o variantă care încalcă o constrângere oprește execuția, în locul unei corecții tacite.}
    \label{fig:motor}
\end{figure}

Două decizii de proiectare necesită explicații suplimentare, deoarece verificările inițiale produceau rezultate aparent corecte, fără a evalua însă condițiile urmărite.

\textbf{Verificarea consistenței se face per rând, nu global.} În prima variantă, eroarea relativă a identităților algebraice era comparată cu o toleranță derivată din valoarea maximă a caracteristicii pe întregul set. Întrucât unele caracteristici au valori maxime de ordinul milioanelor, toleranța rezultată era atât de mare încât verificarea trecea indiferent de conținut. Formularea corectă compară eroarea fiecărui rând cu eroarea aceluiași rând în datele originale: cerința nu este ca identitatea să fie respectată exact, ea nu este respectată nici în datele brute, ci ca perturbarea să nu o degradeze suplimentar.

\textbf{Limitele de domeniu se aplică relativ la rândul original.} Formulate ca praguri absolute derivate din specificațiile protocoalelor, ele respingeau fluxuri prezente în captura reală: setul conține 61 de înregistrări cu valori TTL sub pragul de livrabilitate și o înregistrare a cărei dimensiune medie a pachetului depășește unitatea maximă de transmisie. Deoarece obiectul verificării este perturbarea, nu calitatea datelor de intrare, condiția corectă este ca varianta perturbată să nu încalce o limită mai grav decât o încălca deja rândul original.

Validatorul este supus unei \textbf{testări negative}: șapte variante construite în mod deliberat cu încălcări ale constrângerilor sunt introduse în sistem, cu cerința ca fiecare să fie respinsă. Acestea includ aplicarea umpluturii în sens invers, modificarea unei caracteristici asociate victimei, rescrierea unei caracteristici de identificare și introducerea unor valori nedefinite. Testarea confirmă astfel că validatorul identifică intrările nevalide, nu doar că se execută fără a produce erori.

\section{Măsurarea evaziunii}\label{sec:implmasurare}

Modulul de măsurare încarcă, pentru fiecare model, predicțiile de referință și masca de eligibilitate calculată conform criteriului din subsecțiunea~\ref{subsec:eligibil}, apoi aplică modelul fiecărei variante și clasifică rezultatele conform taxonomiei din subsecțiunea~\ref{subsec:rezultate}.

Variantele nu sunt citite de pe disc, ci regenerate în memorie la momentul măsurării, folosind aceleași generatoare. Generarea fiind deterministă, rezultatul este identic, iar spațiul de stocare necesar scade de la zeci de gigaocteți la câțiva megaocteți. Se persistă doar fișa rulării, rapoartele de validitate și eșantioane de câteva sute de rânduri per variantă, suficiente pentru inspecție manuală.

Pe lângă etichete, modulul înregistrează și distribuțiile de probabilitate, din care calculează deplasarea încrederii modelului: creșterea probabilității atribuite clasei traficului legitim și scăderea probabilității atribuite clasei reale. Aceste mărimi sunt informative chiar și acolo unde eticheta finală nu s-a modificat, deoarece arată cât de aproape de granița de decizie a fost împins fluxul.

Agregarea produce ratele per variantă, defalcarea per clasă de atac, cu marcarea explicită a claselor al căror suport este prea mic pentru o interpretare stabilă, intervalele definite în subsecțiunea~\ref{subsec:interval} pentru perechile de politici, și intensitatea minimă la care fiecare flux evadează pentru prima dată. Ultima mărime răspunde întrebării cât de mare trebuie să fie perturbarea, distinctă de întrebarea câte fluxuri evadează.

Înaintea oricărei agregări se verifică invariantul identității: rata de evaziune pe varianta identică trebuie să fie exact nulă pentru toate modelele. Verificarea este obligatorie și blochează execuția în caz de eșec. Măsurarea pe o singură variantă și încadrarea fiecărui flux în cele trei rezultate posibile sunt reproduse în secțiunile~\ref{sec:anexa-masurare} și \ref{sec:anexa-outcomes}.

\section{Analiza de sensibilitate}\label{sec:implsensibilitate}

Modulul măsoară cât contează fiecare caracteristică pentru decizia modelului, urmând procedura descrisă în subsecțiunea~\ref{subsec:agnostic}, cu cinci repetări pentru fiecare caracteristică, pe modelele înghețate. Metrica principală este scăderea ratei de detecție, din motivele de comensurabilitate expuse în subsecțiunea~\ref{subsec:comensurabil}. Se raportează suplimentar și scăderea de macro-$F1$. Bucla de permutare este reprodusă în secțiunea~\ref{sec:anexa-importanta-cod}.

Analiza restrânsă la o singură clasă de atac a necesitat o corecție de implementare. Permutarea uzuală amestecă valorile coloanei în interiorul eșantionului analizat. Dacă acest eșantion conține o singură clasă, iar caracteristica are aproape aceeași valoare pentru toate observațiile ei, amestecarea nu schimbă practic nimic, iar importanța măsurată apare fals ca fiind nulă, deși caracteristica poate fi tocmai cea pe care se sprijină decizia. Soluția implementată este substituția valorilor dintr-o distribuție \emph{donoare} externă, extrasă din întregul set de testare, care rupe efectiv legătura dintre caracteristică și etichetă.

Un al doilea modul pune în corespondență importanța cu evaziunea. Pentru fiecare transformare se însumează importanțele caracteristicilor pe care le afectează și se corelează seria rezultată cu ratele de evaziune corespunzătoare, prin coeficientul de corelație a rangurilor. Tot aici se calculează fracțiunea accesibilă definită de subsecțiunea~\ref{subsec:accesibila}, prin însumarea importanțelor pe categoriile din tabelul~\ref{tab:taxonomie}.

Un al treilea modul analizează destinațiile rezultatului de tip (c): pentru fiecare model se construiește distribuția claselor către care migrează fluxurile perturbate care rămân semnalate, și se compară profilul de încredere al fluxurilor care evadează cu al celor care nu evadează.

\section{Antrenarea adversarială}\label{sec:impladversarial}

Etapa de antrenare adversarială este singura care produce modele noi. Din acest motiv, ea respectă o regulă suplimentară față de celelalte: nu scrie niciodată în locațiile ocupate de modelele de referință și nici în artefactele derivate din ele. Modelele rezultate și rezultatele lor sunt păstrate în directoare separate, astfel încât toate măsurătorile anterioare rămân valabile și reproductibile după rulare.

\subsection{Reutilizarea mecanismului de perturbare}\label{subsec:reutilizare}

Cerința formulată pentru interfața interactivă în secțiunea~\ref{sec:implinterfata} se aplică identic aici, din același motiv: generatoarele nu sunt reimplementate, ci apelate. Datele de antrenare augmentate se obțin prin aceleași funcții care produc variantele măsurate, cu aceeași propagare a caracteristicilor derivate și aceeași restaurare a tipurilor. Dacă etapa aceasta ar fi construit propriile transformări, modelul ar fi fost antrenat pe o familie de perturbări ușor diferită de cea pe care este ulterior evaluat, iar rezultatul nu ar mai fi spus nimic despre robustețea la perturbările măsurate.

Diferă însă un singur parametru al apelului, cel discutat în subsecțiunea~\ref{subsec:provenienta}: contextul din care generatoarele își extrag statisticile auxiliare este construit exclusiv din partiția de antrenare. Funcția care construiește acest context înregistrează numărul de rânduri al fiecărui cadru sursă primit, iar etapa de augmentare verifică explicit, înainte de a genera orice, că a primit un singur cadru și că acesta are dimensiunea partiției de antrenare. Verificarea este o aserțiune care oprește execuția, nu un avertisment, pentru că scurgerea pe care o previne ar fi invizibilă în rezultate.

\subsection{Eșantionarea variantelor}\label{subsec:esantionare}

Fiecărui flux de atac care primește o copie i se atribuie o variantă prin extragere uniformă peste tipurile admise, urmată de o extragere uniformă peste nivelurile tipului ales. Tipurile marcate nerealizabile sunt excluse din start, la fel și cele scoase deliberat pentru grupul de generalizare.

Implementarea este vectorizată: se extrag simultan atribuirile tuturor rândurilor, apoi acestea se grupează pe perechi de tip și nivel, iar generatorul este apelat o singură dată pentru fiecare grup, pe sub-cadrul corespunzător. Alternativa evidentă, un apel per rând, ar fi fost inutilizabilă, generatoarele operează pe cadre întregi, iar costul fix al unui apel ar fi dominat complet. În forma grupată, augmentarea întregului set de atac durează sub o secundă, ceea ce face practicabilă reeșantionarea la fiecare epocă.

Ordinea rândurilor este fixată înaintea extragerii, iar generatorul de numere aleatoare este inițializat din perechea formată de \textit{seed}-ul rulării și numărul epocii. Varianta atribuită unui flux la o epocă dată se reproduce prin urmare exact, fără a depinde de ordinea în care sunt parcurse loturile sau de planificarea firelor de execuție. Alegerea de a extrage toate atribuirile înaintea construirii loturilor, în locul extragerii la momentul accesării fiecărei observații, este tocmai ceea ce face posibilă această garanție.

\subsection{Două regimuri de augmentare}\label{subsec:regimuri}

Cele două familii de modele primesc datele augmentate în moduri diferite, impuse de mecanismele lor de antrenare.

Ansamblurile de arbori sunt construite dintr-o singură trecere peste un set fix. Ele primesc, prin urmare, setul augmentat o singură dată, cu o variantă atribuită definitiv fiecărui flux de atac. Rețeaua neuronală parcurge însă zeci de treceri peste date, ceea ce permite ca aceleași fluxuri să apară la fiecare trecere cu o altă variantă. Figura~\ref{fig:augmentare} prezintă acest al doilea regim.

Al doilea regim este obținut prin trei metode scurte, apelate câte o dată pe
epocă:

\begin{verbatim}
def epoch_seed(self, epoch: int) -> int:
    """Seed-ul extragerii pentru o epoca."""
    return int(self.seed) * 1_000_003 + int(epoch)

def variants_for_epoch(self, epoch: int) -> pd.DataFrame:
    """Variantele atribuite randurilor de atac la epoca data."""
    frames = []
    for copy_index in range(self.k):
        frames.append(augment.assign_variants(
            len(self.attack_X),
            self.epoch_seed(epoch) + copy_index * 7919,
            self.holdout))
    return pd.concat(frames, ignore_index=True)

def loader_for_epoch(self, epoch: int):
    """DataLoader-ul epocii, amestecat determinist."""
    num, cat, y = self.tensors_for_epoch(epoch)
    return training.make_loader(num, cat, y,
                                config.TRANSFORMER_BATCH_SIZE,
                                shuffle=True,
                                seed=self.epoch_seed(epoch))
\end{verbatim}

Numărul epocii intră în \texttt{epoch\_seed} alături de \textit{seed}-ul rulării,
astfel încât fiecare epocă primește o extragere proprie, dar complet
reproductibilă. Aceeași valoare inițializează și amestecarea loturilor, deci
epoca se reface identic indiferent de ordinea în care sunt parcurse loturile
sau de planificarea firelor de execuție.

Metoda \texttt{loader\_for\_epoch} este singurul punct de legătură cu bucla de
antrenare: ea se transmite acesteia ca argument, iar bucla o apelează la
începutul fiecărei epoci în locul unui set fix. Regimul static se obține
transmițând în locul ei valoarea nulă, caz în care bucla folosește setul primit
o singură dată. Diferența dintre cele două regimuri se reduce, prin urmare, la
un singur argument, iar restul procedurii de antrenare rămâne neschimbat.

\begin{figure}[!ht]
    \centering
    \begin{tikzpicture}[node distance=9mm and 15mm]
        \node[blocmare] (baza) {Tensorii de intrare\\ corespunzători fluxurilor originale};
        \node[blocmare, below=of baza] (esant) {Eșantionarea variantelor\\ (\textit{seed}-ul epocii)};
        \node[blocmare, right=of esant] (gen) {Generare, propagare,\\ preprocesare};
        \node[blocmare, right=24mm of baza] (lot) {Loturile\\ epocii};
        \node[bloc, right=of lot, minimum width=26mm] (ep) {O epocă de\\ antrenare};

        \draw[sageata] (baza) -- (lot);
        \draw[sageata] (esant) -- (gen);
        \draw[sageata] (gen.north) -- (lot.south);
        \draw[sageata] (lot) -- (ep);
        \draw[punctat] (ep.south) -- ++(0,-32mm) -|
            node[eticheta, pos=0.5, below] {epoca următoare} (esant.south);
    \end{tikzpicture}
\caption{Augmentarea dinamică folosită la antrenarea rețelei. Datele fluxurilor originale sunt preprocesate o singură dată, iar la fiecare epocă sunt reeșantionate variantele fluxurilor de atac și este regenerată partea perturbată a loturilor de antrenare.}    \label{fig:augmentare}
\end{figure}

Diferența are o consecință care trebuie raportată, nu ascunsă: rețeaua vede, pe parcursul antrenării, mai multe variante distincte ale aceluiași flux decât arborii. Tentația de a egaliza expunerea, mărind numărul de copii primite de arbori, a fost respinsă deliberat. Numărul de rânduri crește liniar cu acest număr, iar la valori care ar apropia expunerea arborilor de cea a rețelei proporția dintre trafic normal și trafic malițios s-ar deforma sever, iar memoria necesară construirii pădurii ar deveni ea însăși o constrângere. Soluția adoptată este menținerea unui număr mic și mărginit de copii pentru arbori și raportarea explicită a diferenței de expunere, comparația între cele două familii fiind tratată ca orientativă, nu ca paritate controlată.

\subsection{Transmiterea explicită a ponderilor de clasă}\label{subsec:implponderi}

Principiul enunțat în subsecțiunea~\ref{subsec:ponderi} se traduce diferit pentru fiecare model. Random Forest primește în mod obișnuit o indicație simbolică prin care își calculează singur ponderile din frecvențele setului primit. Aici, în locul indicației, i se transmite dicționarul de valori calculat pe setul original. Transformer primește vectorul corespunzător, ordonat conform codificării etichetelor. Modelul cu XGBoost nu folosește ponderi de clasă nici în varianta de referință, deci pentru el nu există nimic de fixat.

\begin{verbatim}
def balanced_weights(y_original) -> dict:
    """Ponderile calculate pe etichetele ORIGINALE."""
    classes = np.unique(np.asarray(y_original))
    values = compute_class_weight("balanced", classes=classes,
                                  y=np.asarray(y_original))
    return {cls: float(w) for cls, w in zip(classes, values)}


def build_rf(numeric_cols, class_weight: dict,
             seed: int) -> Pipeline:
    """class_weight primeste dictionarul inghetat in locul
    sirului "balanced", ca ponderile sa nu se recalculeze pe
    numaratorile augmentate.
    """
    return Pipeline([
        ("rare", RareCategoryGrouper(
            columns=config.CATEGORICAL_COLUMNS,
            min_freq=config.MIN_CATEGORY_FREQUENCY)),
        ("prep", build_preprocessor(numeric_cols)),
        ("clf", RandomForestClassifier(
            n_estimators=config.RF_N_ESTIMATORS,
            class_weight=class_weight,
            n_jobs=config.RF_N_JOBS,
            random_state=seed)),
    ])
\end{verbatim}

Toată diferența față de varianta de referință se află într-un singur argument.
Acolo unde modelul inițial primea șirul \texttt{"balanced"}, prin care
biblioteca își calcula singură ponderile din numărătorile setului primit, aici
primește dicționarul produs de \texttt{balanced\_weights} pe etichetele
originale. Restul componentelor rămân identice, ceea ce face verificabilă
afirmația că grupurile diferă doar prin date.

Ponderile nu sunt calculate o singură dată pentru toate modelele, ci o dată pentru fiecare familie, pe exact mulțimea de etichete pe care varianta de referință le-a calculat. Arborii sunt antrenați pe întreaga partiție de antrenare, iar Transformer pe sub-setul rămas după decuparea validării. Cele două mulțimi diferă, deci și ponderile diferă ușor. Folosirea unui singur vector pentru ambele ar fi introdus tocmai diferența de obiectiv pe care procedura urmărește să o elimine.

\subsection{Decuparea validării înaintea augmentării}\label{subsec:implsplit}

Pentru Transformer, împărțirea în sub-set de antrenare și set de validare se face pe rândurile originale, înaintea oricărei augmentări, cu același \textit{seed} și aceeași fracțiune ca la varianta de referință. Egalitatea cu împărțirea originală este verificată prin compararea mulțimilor de identificatori, nu presupusă.

Copiile perturbate se generează apoi exclusiv din rândurile aflate de partea de antrenare. O verificare separată confirmă că mulțimea identificatorilor sursă ai rândurilor generate este inclusă în cea a sub-setului de antrenare și disjunctă de cea a validării. Fără această condiție, o copie a unui flux de antrenare ar putea ajunge în validare, iar scorul care decide oprirea antrenării ar fi calculat parțial pe date derivate din antrenare.

Setul de validare rămâne integral nemodificat. Criteriul de selecție a modelului final este cel din varianta de referință, calculat pe trafic curat. Alegerea este deliberată și este argumentată în subsecțiunea~\ref{subsec:confuzii}: dacă oprirea ar fi condusă de un scor măsurat pe date perturbate, grupul apărat ar diferi de cel de referință atât prin date, cât și prin criteriul de selecție, iar comparația nu ar mai fi controlată. Robustețea se măsoară după antrenare, nu în timpul ei.

\subsection{Reluarea și artefactele de reproductibilitate}\label{subsec:implreluare}

Antrenarea rețelei de referință salvează periodic starea completă, inclusiv starea generatoarelor de numere aleatoare și pe cea a mecanismului care produce loturile, ceea ce permite reluarea exactă după o întrerupere. Pentru grupurile augmentate, mecanismul nu mai este suficient: loturile fiind regenerate la fiecare epocă, starea salvată nu descrie complet situația de reluat. Soluția adoptată este reluarea la granularitatea grupului: fiecare model deja salvat este sărit la o rulare ulterioară, iar o întrerupere costă cel mult grupul aflat în curs. Aceeași regulă se aplică și etapei de evaluare, unde fiecare rulare deja măsurată este sărită.

Seturile de date augmentate nu sunt salvate integral, deoarece pot fi reconstruite în mod determinist într-un timp redus. Persistarea celor patru seturi, fiecare conținând câteva sute de mii de rânduri, ar crește inutil necesarul de stocare și de memorie. Pentru fiecare set este salvată în schimb o amprentă criptografică, calculată după ordonarea coloanelor reprezentării tabelare. Aceasta permite verificarea faptului că datele reconstruite ulterior sunt identice cu cele utilizate inițial.

Pentru fiecare rulare este generată o fișă care înregistrează parametrii necesari reproducerii experimentului: \textit{seed}-ul folosit, numărul de copii generate pentru fiecare flux, variantele permise, probabilitățile de eșantionare pentru fiecare tip și variantă, ponderile utilizate, precum și compoziția și amprenta fiecărui set de date. Manifestul include, de asemenea, proveniența contextelor de generare, rezultatele verificărilor de integritate și pragurile stabilite pentru criteriile de evaluare. Valorile efectiv folosite sunt listate în secțiunea~\ref{sec:anexa-config}. Aceste praguri sunt înregistrate înainte de începerea rulării, conform procedurii descrise în subsecțiunea~\ref{subsec:preinregistrare}.

\subsection{Evaluarea pe mulțimea comună}\label{subsec:implcohorta}

Modulul de evaluare parcurge, pentru fiecare model antrenat, aceeași grilă de variante ca etapa de măsurare sistematică și înregistrează rezultatul fiecărui flux la fiecare variantă într-o matrice de coduri întregi pe un octet. Alegerea de a reține rezultatul tuturor fluxurilor de atac, nu doar al celor semnalate de modelul respectiv, este ceea ce permite recitirea aceleiași măsurători atât pe mulțimea proprie de fluxuri semnalate, cât și pe mulțimea comună definită în subsecțiunea~\ref{subsec:cohorta}, fără reluarea predicțiilor. Întreaga colecție de matrice ocupă câteva zeci de megaocteți.

Mulțimea comună se obține prin intersecția măștilor de eligibilitate ale tuturor rulărilor aceleiași familii de modele, citite din rândul corespunzător variantei identice al fiecărei matrice. Pe mulțimea astfel fixată se verifică din nou invariantul identității, iar apoi se calculează ratele globale, defalcarea per clasă de atac și mediile peste repetări.

Variantele pe care se face această măsurare sunt generate cu contextul complet, identic cu cel folosit la măsurarea sistematică, pentru ca perturbările evaluate să fie aceleași rând cu rând. Restricția la partiția de antrenare se aplică exclusiv generării datelor de antrenare. În momentul evaluării, modelele sunt deja antrenate și înghețate, deci nu mai pot fi influențate de proveniența statisticilor. Cele două roluri ale aceluiași mecanism sunt separate explicit în cod, prin două obiecte de context distincte, construite o singură dată și transmise mai departe.


\section{Interfața interactivă}\label{sec:implinterfata}

Interfața este realizată ca aplicație web locală și expune trei vederi: rezultatul agregat pe un grup de fluxuri, inspecția unui flux individual și tabelele și figurile măsurate sistematic.

Interfața nu își construiește propriile perturbări. Ea apelează exact aceleași funcții folosite la măsurarea sistematică, singura diferență fiind că valorile parametrilor vin de la utilizator, nu dintr-o listă stabilită dinainte. Cerința aceasta, formulată în secțiunea~\ref{sec:structura}, are un motiv practic. Dacă interfața ar avea propriul cod de perturbare, cele două implementări ar putea începe să difere după o modificare făcută doar în una dintre ele, iar cifrele afișate pe ecran nu ar mai corespunde celor raportate în lucrare.

Controalele sunt grupate într-un formular, astfel încât nicio recalculare să nu se producă până la acționarea butonului de rulare. Rațiunea este că valorile afișate trebuie să corespundă întotdeauna ultimei rulări, nu câmpurilor editate între timp. În absența acestei separări, un utilizator care modifică un parametru fără a-l aplica ar vedea rezultate care nu îi corespund. Parametrii efectiv utilizați sunt afișați într-un rezumat, tocmai pentru ca discrepanța dintre câmpurile editate și rezultatele afișate să fie explicită.

Modelele și datele se încarcă o singură dată și sunt păstrate în memorie între interacțiuni. Selecția fluxurilor dintr-o clasă se face prin eșantionare aleatoare cu \textit{seed} fix, nu prin preluarea primelor înregistrări: setul fiind ordonat după identificator, primele înregistrări nu sunt reprezentative pentru distribuția clasei, iar ratele afișate ar fi fost sistematic deplasate.

Dacă o combinație de parametri produce o variantă care încalcă o constrângere, interfața o semnalează explicit în locul afișării unui rezultat. Fără această verificare, utilizatorul ar putea obține o rată de evaziune ridicată pentru un flux imposibil fizic, iar valoarea ar fi lipsită de sens.

\section{Structurile de date persistate}\label{sec:structuri}

Comunicarea între etape se face exclusiv prin fișiere tabelare, ceea ce permite atât reluarea independentă a etapelor, cât și inspecția manuală a rezultatelor intermediare. Structurile principale sunt descrise în tabelul~\ref{tab:structuri}.

\begin{table}[H]
    \caption{Structurile de date schimbate între etape}
    \centering
    \begin{tabular}{|l|p{7.4cm}|}
        \hline
        \textbf{Structură} & \textbf{Conținut} \\
        \hline
        Predicții de referință & Identificator, etichetă reală și predicția fiecărui model pe traficul nemodificat \\
        \hline
        Distribuții de probabilitate & Matricea probabilităților per clasă, per model \\
        \hline
        Indicatori de eligibilitate & Masca definită în subsecțiunea~\ref{subsec:eligibil}, per model \\
        \hline
        Fișa variantelor & Câte o linie per variantă: tip, nivel, parametri, caracteristici modificate, statistici de rezolvare a contorului, cazuri degenerate, verificări executate și eșuate \\
        \hline
        Raportul de validitate & Câte o linie per verificare și per variantă, cu numărul de încălcări \\
        \hline
        Rezultate per flux & Identificator, etichetă reală, predicție și categorie de rezultat, pentru fiecare variantă și model \\
        \hline
        Agregări & Rate per variantă, defalcări per clasă, intervale, intensități minime \\
        \hline
    \end{tabular}
    \label{tab:structuri}
\end{table}

Fișa rulării are un rol care depășește evidența: el conține toate mărimile necesare pentru a stabili \emph{cât} din rezultatul unei variante depinde de decizii de rezervă. Numărul de rânduri rezolvate prin fiecare dintre cele trei niveluri descrise în subsecțiunea~\ref{subsec:ct} și numărul de cazuri degenerate din propagare sunt înregistrate explicit, astfel încât un rezultat obținut majoritar prin mecanisme de rezervă să poată fi identificat ca atare.

Rezultatele per flux sunt salvate comprimat, întrucât ele reprezintă produsul cartezian dintre fluxurile eligibile și variante. Ele sunt necesare analizei de sensibilitate, care compară profilul fluxurilor evadate cu al celor neevadate, comparație imposibilă pornind exclusiv de la agregări.

\section{Jurnalizarea și tratarea erorilor}\label{sec:jurnalizare}

Fiecare etapă produce un jurnal de execuție, scris simultan în consolă și într-un fișier din directorul de rezultate. Jurnalul înregistrează parametrii efectiv utilizați, dimensiunile datelor prelucrate, rezultatul fiecărei verificări și duratele etapelor intermediare.

Principiul adoptat în tratarea erorilor este eșecul explicit și timpuriu. O condiție încălcată oprește execuția cu un mesaj care indică verificarea și contextul, în locul continuării cu o valoare implicită. Alegerea este justificată de natura experimentală a codului: o etapă care continuă după o încălcare produce rezultate care par valide și care ar fi interpretate ca atare.

Continuarea procesului este condiționată de trecerea a trei verificări: validarea variantelor perturbate, concordanța dintre predicțiile modelului reîncărcat și cele ale modelului păstrat în memorie, precum și verificarea rezultatelor pentru varianta nemodificată. Dacă una dintre aceste verificări eșuează, etapele următoare nu mai sunt executate.


\section{Asigurarea reproductibilității}\label{sec:reproductibilitate}

Reproductibilitatea este esențială pentru interpretarea corectă a rezultatelor. Experimentele au fost executate în momente diferite, uneori la intervale de câteva săptămâni, astfel încât rezultatele pot fi comparate numai dacă repetarea unei rulări în aceleași condiții conduce la aceleași valori.

Pentru toate operațiile aleatoare este utilizată aceeași valoare inițială, definită în configurație. Aceasta este transmisă bibliotecilor folosite pentru prelucrarea datelor și antrenarea modelelor, inclusiv operațiilor executate pe placa grafică. În plus, inferența se desfășoară pe un singur fir de execuție, din motivele prezentate în subsecțiunea~\ref{subsec:inferenta}, iar pentru antrenarea rețelei sunt selectate operații deterministe, chiar dacă acestea presupun un timp de execuție mai mare.

La fiecare rulare completă se salvează un fișier cu hiperparametrii folosiți, versiunile bibliotecilor, ordinea caracteristicilor așteptată de model și un rezumat al setului de date. Utilitatea lui s-a văzut într-un caz concret. Instalarea unei biblioteci de analiză a coborât tacit versiunea unei dependențe numerice, iar rezultate deja obținute s-au schimbat. Fără versiunile înregistrate, diferența ar fi fost pusă pe seama unei modificări de cod și căutată acolo unde nu era.

Două verificări încheie această parte. Fiecare model salvat este reîncărcat de pe disc și pus să prezică din nou, iar predicțiile trebuie să coincidă cu cele obținute înainte de salvare. Apoi varianta neperturbată produsă de motorul de perturbare este trecută prin toate cele trei modele și trebuie să reproducă predicțiile de referință rând cu rând. Împreună cu invariantul identității, ele garantează că diferențele din capitolul următor provin din perturbări, nu din variații ale prelucrării.
